\documentclass[11pt]{article}

% --- packages ---
\usepackage[margin=1in]{geometry}
\usepackage{mathptmx}            % times-like serif, paper look
\usepackage{graphicx}
\usepackage{booktabs}
\usepackage{amsmath}
\usepackage{caption}
\usepackage{subcaption}
\usepackage{xcolor}
\usepackage[colorlinks=true,linkcolor=black,citecolor=black,urlcolor=blue]{hyperref}

\captionsetup{font=small,labelfont=bf}
\setlength{\parskip}{0.35em}

% --- title ---
\title{\vspace{-1.5em}\textbf{Frozen Vision--Language Features for\\ Fine-Grained 100-Class Transfer Learning}\vspace{-0.3em}}
\author{
  Jay Jani (CruzID: \texttt{XXXXX})\\
  % add up to two more members here
  \small CSE 144: Applied Machine Learning --- Final Project, UC Santa Cruz
}
\date{Spring 2026}

\begin{document}
\maketitle

\begin{abstract}
\noindent
I tackle a 100-class fine-grained image classification task spanning four visual
domains (food, flowers, cars, and aircraft) with only about ten labeled images per
class. With so little data, fully fine-tuning a large network overfits. Instead I
use transfer learning through \emph{frozen} feature extraction: I keep two
pretrained SigLIP-2 vision--language backbones fixed and train only lightweight
linear classifiers on top of their embeddings, combined with zero-shot text
classifiers built from the class names. A small ensemble of these heads, with
weights tuned by cross-validation, reaches \textbf{96.3\%} five-fold accuracy and
\textbf{98.1\%} on the Kaggle public leaderboard. I show that this frozen approach
beats my best fine-tuning runs while being far cheaper and fully reproducible.
\end{abstract}

% ======================================================================
\section{Introduction}

The goal is to classify images into 100 categories sampled from several different
source datasets, so the classes are \emph{fine-grained}: many cars, several flower
species, dishes, and aircraft models that differ only in small details. The
training set has roughly ten images per class, which is far too few to train a deep
network from scratch.

Transfer learning is the natural fit. A model pretrained on a very large image
corpus already encodes general visual features; I only need to adapt those
features to the 100 labels. The usual recipe is to fine-tune the pretrained
weights, but with $\sim$10 images per class fine-tuning quickly memorizes the
training set. I therefore take the opposite end of the transfer-learning
spectrum: I \emph{freeze} the backbone entirely and learn only a small classifier
on top of its fixed embeddings.

My approach combines (i) linear probes on frozen SigLIP-2 embeddings, (ii)
zero-shot text classifiers that compare each image to the encoded class names, and
(iii) a cross-validated weighted ensemble of these heads. The final system reaches
96.3\% accuracy under five-fold cross-validation and 98.1\% on the public
leaderboard, comfortably above the 60\% baseline, while training in seconds on CPU
once the embeddings are cached.

% ======================================================================
\section{Dataset}

\paragraph{Size and structure.} The data contain 100 classes drawn from four
domains. The training set has 1{,}079 images organized as \texttt{train/<label>/*.jpg},
and the test set has 1{,}036 unlabeled images in \texttt{test/*.jpg}. Folder names
are the integer labels, so class ``\texttt{0}'' maps to label $0$, class
``\texttt{1}'' to label $1$, and so on; I keep this mapping exactly to match the
grader.

\paragraph{Domains.} The 100 classes are contiguous blocks of 25: labels 0--24 are
food, 25--49 are flowers, 50--74 are cars, and 75--99 are aircraft. I use this
structure later to build domain-aware text prompts.

\paragraph{Preprocessing.} Images come in varied sizes, so each is resized to the
backbone's native resolution (384$\times$384 for one model, 512$\times$512 for the
other) and normalized with that model's mean and standard deviation. At inference
I use simple test-time augmentation (TTA): four views per image (a center crop, a
square resize, and their horizontal flips), whose embeddings are averaged. I do no
random training-time augmentation, because the backbone is frozen and the probe is
trained on fixed feature vectors.

% ======================================================================
\section{Method}

\subsection{Frozen backbones}
I use two pretrained SigLIP-2 vision--language models from the
\texttt{open\_clip} library: a ViT-g/16 at 384px (1536-d embeddings) and a
SO400M/16 at 512px (1152-d embeddings). Both are trained with a contrastive
image--text objective on web-scale data. I chose contrastive
vision--language models on purpose: the classes are nameable, web-common concepts,
which is exactly what these models align well. Crucially, the backbones are never
updated; I only run them forward to extract embeddings, which I cache to disk.

\subsection{Linear probes}
For each backbone I train a multinomial logistic-regression probe on the
L2-normalized embeddings (\texttt{scikit-learn}, inverse regularization $C=10$,
balanced class weights). I also train a third probe on the \emph{concatenation} of
both backbones' embeddings (early fusion), which captures complementary information
the two models encode differently.

\subsection{Zero-shot text heads}
Each backbone also has a text encoder. For every class I encode its name with a
few short, domain-specific prompt templates (for example
``\emph{a photo of the car model \{\}}'' for cars, or
``\emph{a close-up photo of a \{\} flower}'' for flowers) and average the resulting
text vectors. An image is then scored by cosine similarity to each class vector.
I calibrate the softmax temperature of each text head by minimizing
cross-entropy on the training folds, which sharpens the probabilities and improves
accuracy at no extra cost.

\subsection{Cross-validated ensemble}
The above gives five candidate heads: two linear probes, two text heads, and the
concatenated probe. I combine their predicted probabilities with non-negative
weights and pick the argmax. The weights are tuned to maximize five-fold
out-of-fold (OOF) accuracy using random-restart coordinate ascent, which avoids
the local optima a plain greedy search falls into. Weak heads are free to receive
weight zero, so the search effectively performs model selection.

% ======================================================================
\section{Experiments}

\paragraph{Validation protocol.} All design choices are made on a leak-free
5-fold stratified cross-validation of the training set. Each test prediction comes
from a probe trained only on the training data; no test label is ever used.

\paragraph{Which backbone?} I compared several strong pretrained backbones by the
accuracy of a single linear probe on their frozen features
(Table~\ref{tab:backbones}). The two SigLIP-2 models are clearly the strongest. I
also tried adding the others to the ensemble; in every case the cross-validated
weight search drove them to zero, meaning they added no information beyond the two
SigLIP-2 models. Self-supervised features (DINOv2) lag the contrastive
vision--language models, consistent with the classes being nameable concepts.

\begin{table}[t]
\centering
\caption{Single linear-probe accuracy (5-fold OOF) for different frozen backbones.
Only the two SigLIP-2 models receive non-zero weight in the final ensemble.}
\label{tab:backbones}
\begin{tabular}{llc}
\toprule
Backbone & Pretraining & Probe OOF \\
\midrule
SigLIP-2 ViT-g/16 (384) & contrastive (WebLI) & \textbf{0.944} \\
SigLIP-2 SO400M/16 (512) & contrastive (WebLI) & \textbf{0.942} \\
DFN ViT-H/14 (378) & contrastive & 0.932 \\
ConvNeXt-XXL & contrastive (LAION-2B) & 0.913 \\
EVA02-L/14 (336) & contrastive & 0.906 \\
MetaCLIP-2 ViT-H/14 & contrastive & 0.894 \\
DINOv2 ViT-g/14 & self-supervised & 0.846 \\
\bottomrule
\end{tabular}
\end{table}

\paragraph{Frozen probing vs.\ fine-tuning.} I also tried fully fine-tuning a
backbone with layer-wise learning-rate decay, label smoothing, and stochastic
depth. Despite this regularization, the small training set caused overfitting and
my best fine-tuning result reached only about 95.4\%, below the frozen ensemble
and far more expensive and sensitive to seeds. This is the central reason my final
system is frozen: with $\sim$10 images per class, adapting only a linear head on
strong fixed features generalizes better than moving the backbone weights.

\paragraph{Probe hyperparameters.} The probe accuracy is flat across $C \in [3,
100]$; an $\ell_2$ penalty with the \texttt{lbfgs} solver beats $\ell_1$, and
balanced class weights match the default at the plateau. I therefore fix $C=10$.
A nonlinear probe (MLP) did not help, because the contrastive features are already
close to linearly separable.

% ======================================================================
\section{Results}

\paragraph{Member and ensemble accuracy.} Table~\ref{tab:members} lists the
accuracy of each ensemble member and of the final weighted combination. The
ensemble reaches 96.3\% OOF accuracy, about two points above the best single head,
and the per-fold standard deviation is small ($\pm 0.008$).

\begin{table}[t]
\centering
\caption{Final ensemble members and combined accuracy (5-fold OOF).}
\label{tab:members}
\begin{tabular}{lc}
\toprule
Head & OOF accuracy \\
\midrule
SigLIP-2 ViT-g probe        & 0.944 \\
SigLIP-2 ViT-g text         & 0.951 \\
SigLIP-2 SO400M probe       & 0.942 \\
SigLIP-2 SO400M text        & 0.943 \\
Concatenated probe          & 0.948 \\
\midrule
\textbf{Weighted ensemble}  & \textbf{0.963} \\
\bottomrule
\end{tabular}
\end{table}

\paragraph{Probe convergence.} Figure~\ref{fig:curve} shows training and
validation loss and accuracy for a linear head trained on the frozen ViT-g
embeddings, across all five folds. The probe converges quickly and the small
gap between train and validation curves confirms it is not overfitting, which is
expected when only a linear layer is learned.

\begin{figure}[t]
\centering
\includegraphics[width=0.92\linewidth]{probe_curve.png}
\caption{Five-fold loss and accuracy curves for a linear probe on frozen SigLIP-2
ViT-g embeddings. Faint lines are individual folds; bold lines are the mean. Final
validation accuracy is $0.946 \pm 0.004$.}
\label{fig:curve}
\end{figure}

\paragraph{Leaderboard.} My final submission scores \textbf{98.1\%} on the Kaggle
public leaderboard. Because the public split is only about 10\% of the test set
($\sim$114 images), a single image is worth nearly one percent, so this number is
noisy; I treat the 96.3\% cross-validation estimate as the more reliable measure
of true performance. Both are well above the 60\% baseline.

% ======================================================================
\section{Discussion}

\paragraph{What worked.} The biggest win was choosing transfer by frozen feature
extraction rather than fine-tuning. Strong contrastive backbones already separate
these classes almost linearly, so a tiny classifier on fixed features both
generalizes better and is trivially reproducible. Temperature-calibrated,
domain-aware text prompts and the concatenated early-fusion probe each gave small
additional gains.

\paragraph{The ensemble is decorrelation-limited.} Once both SigLIP-2 models are
in the ensemble, adding more accurate or more backbones barely moves the combined
score. The reason is that ensembling helps only when members make
\emph{different} mistakes; the remaining errors are largely shared across all
heads. This explains why every extra backbone I tried was weighted to zero.

\paragraph{Failure cases.} A clear example is two car classes that correspond to
nearly identical vehicle variants: the images are visually indistinguishable, so
no model can separate them reliably from pixels alone. Errors like these are
effectively irreducible and account for part of the remaining $\sim$3.7\%.

\paragraph{Limitations and next steps.} The method depends on having class names
that match web concepts, which is what makes the zero-shot text heads useful; it
would help less on abstract categories. I also explored self-training
(pseudo-labeling confident test images), but its effect flipped sign across random
seeds, i.e.\ it was within noise, so I did not rely on it. A promising direction
is light, heavily regularized fine-tuning of only the last few blocks, which might
combine the stability of frozen features with a little task-specific adaptation.

% ======================================================================
\section{Reproducibility}

\paragraph{Determinism.} I fix a single random seed (42) that seeds Python,
NumPy, and PyTorch, and enable deterministic cuDNN. The probe and cross-validation
are CPU-based and seeded, so results are bit-reproducible on a given machine. I
report 5-fold cross-validated accuracy with its standard deviation rather than a
single split.

\paragraph{Environment.} PyTorch 2.11 (CUDA 12.8), \texttt{open\_clip},
\texttt{scikit-learn} 1.6.1, NumPy 2.0.2. A hardware preset auto-selects batch size
and precision for the detected GPU (T4 / L4 / A100). Backbone embeddings are
computed once and cached, after which the full pipeline runs in CPU-seconds. A cold
run (encoding both backbones from scratch) takes about 9 minutes on an A100 80\,GB
($\sim$67\,GB peak) and about 25 minutes on an L4 24\,GB ($\sim$18\,GB peak); once
the embeddings are cached, re-running the probe and ensemble takes only seconds.

\paragraph{Commands.} The full pipeline (encode, probe, ensemble, write
\texttt{submission.csv}) runs with:
\begin{quote}\small\ttfamily
bash run.sh
\end{quote}
or directly:
\begin{quote}\small\ttfamily
python src/ensemble\_probes.py \textbackslash\\
\hspace*{1em}--backbone ViT-gopt-16-SigLIP2-384:webli:cache/gopt \textbackslash\\
\hspace*{1em}--backbone ViT-SO400M-16-SigLIP2-512:webli:cache/so400m512 \textbackslash\\
\hspace*{1em}--write submission.csv
\end{quote}
The probe curve figure is produced by
\texttt{python src/probe\_curve.py --emb cache/gopt/train.npy --out probe\_curve.png}.

% ======================================================================
\section{Team Contributions}

\begin{itemize}
  \item \textbf{Jay Jani:} method design, backbone and ensemble experiments,
        code, report. % edit/split as needed
  % \item Member 2: ...
  % \item Member 3: ...
\end{itemize}

% ======================================================================
\begin{thebibliography}{9}
\small
\bibitem{siglip} X. Zhai, B. Mustafa, A. Kolesnikov, L. Beyer.
  \emph{Sigmoid Loss for Language Image Pre-Training (SigLIP)}. ICCV, 2023.
\bibitem{siglip2} M. Tschannen et al.
  \emph{SigLIP 2: Multilingual Vision--Language Encoders}. arXiv, 2025.
\bibitem{clip} A. Radford et al.
  \emph{Learning Transferable Visual Models from Natural Language Supervision
  (CLIP)}. ICML, 2021.
\bibitem{openclip} G. Ilharco et al. \emph{OpenCLIP}.
  \url{https://github.com/mlfoundations/open_clip}, 2021.
\bibitem{sklearn} F. Pedregosa et al.
  \emph{Scikit-learn: Machine Learning in Python}. JMLR, 2011.
\bibitem{dinov2} M. Oquab et al.
  \emph{DINOv2: Learning Robust Visual Features without Supervision}. TMLR, 2024.
\end{thebibliography}

\end{document}
